% ch4_answers.tex -- answers to the Quick check questions of Module 5, Chapter 4.
\begin{chapteranswers}{4}{Gravity, Orbits and Kepler's Laws}
  % ---- Section 4.1
  \item $F \propto \dfrac{Mm}{r^2}$. Doubling one mass multiplies $F$ by 2. Tripling $r$ divides
        $F$ by $3^2 = 9$. New force $= F \times \dfrac{2}{9}$.
        Answer: \final{$\dfrac{2F}{9}$} (about $0.22F$).
  \item (a) By Newton's third law the ball pulls the Earth with a force of the same size in the
        opposite direction: \final{\qty{4.9}{N} upwards} (towards the ball).
        (b) The Earth's acceleration is
        $a = \dfrac{F}{m} = \dfrac{\qty{4.9}{N}}{\qty{6.0e24}{kg}} = \qty{8.2e-25}{m.s^{-2}}$.
        The Earth's mass is so huge that this acceleration is far too small to notice.
  \item $g = \dfrac{GM}{r^2} = \dfrac{\num{6.67e-11} \times \num{7.35e22}}{(\num{1.74e6})^2}
        = \dfrac{\num{4.902e12}}{\num{3.028e12}} = \qty{1.619}{N.kg^{-1}}$.
        Answer: \final{$g = \qty{1.62}{N.kg^{-1}}$} (about one sixth of the Earth's value).
  \item $g \propto \dfrac{1}{r^2}$. For $g$ to be $\dfrac19$ of the surface value, $r^2$ must be
        9 times bigger, so $r$ must be $\sqrt 9 = 3$ times bigger: $r = 3r_\mathrm{E}$.
        The altitude is $h = r - r_\mathrm{E} = 3r_\mathrm{E} - r_\mathrm{E} = 2r_\mathrm{E}
        = 2 \times \qty{6.371e6}{m} = \qty{1.274e7}{m}$.
        Answer: \final{$h = \qty{1.27e7}{m}$} (about \qty{12700}{km}).
  \item At the zero point the two fields have the same size and opposite directions. Let $x$ be
        the distance from the Earth's centre. Then the distance from the Moon's centre is $d - x$,
        with $d = \qty{3.84e8}{m}$.
        \begin{align*}
          \frac{GM_\mathrm{E}}{x^2} &= \frac{GM_\mathrm{M}}{(d-x)^2} && \text{the fields have the same size}\\
          \frac{M_\mathrm{E}}{x^2} &= \frac{M_\mathrm{M}}{(d-x)^2} && \text{cancel $G$}\\
          \frac{(d-x)^2}{x^2} &= \frac{M_\mathrm{M}}{M_\mathrm{E}} && \text{multiply by $(d-x)^2$; divide by $M_\mathrm{E}$}\\
          \frac{(d-x)^2}{x^2} &= \frac{\num{7.35e22}}{\num{6.0e24}} = 0.01225 && \text{substitute}\\
          \frac{d-x}{x} &= \sqrt{0.01225} = 0.1107 && \text{square root of both sides}\\
          \frac{d}{x} - 1 &= 0.1107 && \text{split the fraction: $\frac{d-x}{x} = \frac{d}{x} - 1$}\\
          \frac{d}{x} &= 1.1107 && \text{add 1 to both sides}\\
          x &= \frac{\qty{3.84e8}{m}}{1.1107} = \qty{3.457e8}{m} && \text{rearrange for $x$}
        \end{align*}
        Answer: \final{$x = \qty{3.46e8}{m}$} from the Earth's centre (90\,\% of the way to the
        Moon). Check: both fields are \qty{3.35e-3}{N.kg^{-1}} there.
  % ---- Section 4.2
  \item $r = r_\mathrm{E} + h = \qty{6.371e6}{m} + \qty{2.000e6}{m} = \qty{8.371e6}{m}$.
        $v = \sqrt{\dfrac{GM}{r}} = \sqrt{\dfrac{\qty{4.002e14}{m^3.s^{-2}}}{\qty{8.371e6}{m}}}
        = \sqrt{\qty{4.781e7}{m^2.s^{-2}}} = \qty{6.914e3}{m.s^{-1}}$.
        Answer: \final{$v = \qty{6.91e3}{m.s^{-1}}$}.
  \item Speed: $v = \sqrt{\dfrac{GM}{r}}$, so $v \propto \dfrac{1}{\sqrt r}$. With $r$ 4 times
        bigger, $v$ is divided by $\sqrt 4 = 2$: \final{$v_\mathrm{B} = \tfrac12 v_\mathrm{A}$}.
        Period: $T^2 \propto r^3$, so $T \propto r^{3/2}$. With $r$ 4 times bigger, $T$ is
        multiplied by $4^{3/2} = \sqrt{64} = 8$: \final{$T_\mathrm{B} = 8\,T_\mathrm{A}$}.
        (Check: B goes 4 times as far at half the speed, so it takes $4 \times 2 = 8$ times as long.)
  \item Use $T = 2\pi\sqrt{\dfrac{r^3}{GM}}$ with the Earth's $GM$.
        \begin{align*}
          T &= 2\pi\sqrt{\frac{(\qty{3.84e8}{m})^3}{\qty{4.002e14}{m^3.s^{-2}}}}
            = 2\pi\sqrt{\frac{\qty{5.662e25}{m^3}}{\qty{4.002e14}{m^3.s^{-2}}}}\\
          &= 2\pi \times \qty{3.761e5}{s} = \qty{2.363e6}{s}
        \end{align*}
        In days: $\num{2.363e6} \div \num{86400} = 27.4$. Answer: \final{$T = 27.4$ days}
        (the measured value is 27.3 days).
  \item Gravity still acts on the astronaut: $g = \qty{8.7}{m.s^{-2}}$ at that height. The
        astronaut and the ISS are both in free fall around the Earth, with the same acceleration
        ($g$ towards the Earth's centre). Because they fall together, the floor does not push on
        the astronaut: the normal force is zero. We feel our weight through this push, so the
        astronaut feels weightless (apparent weightlessness).
  % ---- Section 4.3
  \item Perigee and apogee are the closest and farthest points, so
        $r_\mathrm{p} v_\mathrm{p} = r_\mathrm{a} v_\mathrm{a}$ can be used.
        $v_\mathrm{a} = \dfrac{r_\mathrm{p} v_\mathrm{p}}{r_\mathrm{a}}
        = \dfrac{\qty{7.0e6}{m} \times \qty{8.73e3}{m.s^{-1}}}{\qty{1.4e7}{m}} = \qty{4.37e3}{m.s^{-1}}$.
        Answer: \final{$v_\mathrm{a} = \qty{4.4e3}{m.s^{-1}}$} (twice as far, so half as fast).
  \item Same central body (the Sun), so $\dfrac{r_1^3}{T_1^2} = \dfrac{r_2^3}{T_2^2}$, which
        gives $T_2^2 = T_1^2 \left(\dfrac{r_2}{r_1}\right)^3$ (as in \cref{ch4:we:cometT}).
        Measure distances in units of the Earth's orbit radius (AU). Earth: $r_1 = 1$,
        $T_1 = \qty{1}{\yr}$. Ceres: $r_2 = 2.77$ in the same units.
        $T_2^2 = (\qty{1}{\yr})^2 \times 2.77^3 = \qty{21.25}{\yr^2}$, so
        $T_2 = \sqrt{\qty{21.25}{\yr^2}} = \qty{4.610}{\yr}$.
        Answer: \final{$T = \qty{4.61}{\yr}$}.
  \item In a circular orbit, gravity always points to the centre, at \ang{90} to the velocity.
        A force at \ang{90} to the motion does no work ($W = Fs\cos 90^\circ = 0$), so the
        kinetic energy and the speed stay constant. In an elliptical orbit, gravity is not at
        \ang{90} to the velocity (except at perihelion and aphelion). When the planet moves
        closer, part of gravity acts along the motion and does positive work, so it speeds up.
        When it moves away, part of gravity acts against the motion, so it slows down.
  \item Only gravity acts on the comet, so its total energy $K + U$ stays constant. Its
        potential energy $U = -\dfrac{GMm}{r}$ is highest (least negative) when $r$ is largest,
        at aphelion. Total energy is constant, so the kinetic energy $K$ is smallest there.
        Smallest kinetic energy means smallest speed, so the comet moves most slowly at aphelion.
  % ---- Section 4.4
  \item (1) The period is one day (\qty{24}{h}, strictly 23\,h\,56\,min). (2) The orbit is a
        circle in the plane of the equator. (3) The satellite moves from west to east, the same
        direction as the Earth spins.
  \item Change the period to seconds, find $r$ with Kepler's third law, then subtract $r_\mathrm{E}$.
        \begin{align*}
          T &= 14.1 \times \qty{3600}{s} = \qty{5.076e4}{s} && \text{hours to seconds}\\
          r^3 &= \frac{GMT^2}{4\pi^2} && \text{Kepler's third law}\\
          &= \frac{\qty{4.002e14}{m^3.s^{-2}} \times (\qty{5.076e4}{s})^2}{39.48}\\
          &= \frac{\qty{1.031e24}{m^3}}{39.48} = \qty{2.612e22}{m^3}\\
          r &= \sqrt[3]{\qty{2.612e22}{m^3}} = \qty{2.967e7}{m} && \text{cube root}\\
          h &= r - r_\mathrm{E} = \qty{2.967e7}{m} - \qty{0.6371e7}{m} && \text{same power of ten}\\
          &= \qty{2.330e7}{m}
        \end{align*}
        Answer: \final{$h = \qty{2.33e7}{m}$} (about \qty{23300}{km}).
  \item No. Its period matches the Earth's spin, so it is geosynchronous. But the orbit is not in
        the plane of the equator. Seen from the ground, the satellite moves north and south of the
        equator each day, so it does not stay above one point. A geostationary orbit must be
        above the equator.
  \item A low Earth orbit is only a few hundred kilometres up, while a geostationary orbit is
        about \qty{36000}{km} up, nearly 100 times further away. A camera sees much finer detail
        from close up, so small objects such as buildings can be picked out. The satellite moves
        over many places each day, and a polar orbit lets it photograph the whole Earth.
\end{chapteranswers}
